% =====================================================================
%  Nguồn vector cho Hình 4.16 — Kiến trúc pipeline dịch vụ OCR (không PDF text-layer)
%  Khớp ĐÚNG code app.py: 3 engine tự chủ (vietocr / paddleocr / paddledet_viet).
%  Biên dịch: xelatex src_img27.tex -> src_img27.pdf -> copy sang img27.pdf
% =====================================================================
\documentclass[border=10pt]{standalone}
\usepackage{fontspec}
\setmainfont{Arial}
\usepackage{tikz}
\usetikzlibrary{arrows.meta,shapes.geometric,positioning,calc}

\tikzset{
  io/.style   ={rounded corners=8pt, draw=gray!60!black, fill=gray!8, line width=0.8pt,
                align=center, inner sep=6pt, font=\small, minimum height=1cm},
  proc/.style ={draw=gray!60!black, fill=gray!5, line width=0.8pt,
                align=center, inner sep=6pt, font=\small, minimum height=1cm},
  dec/.style  ={diamond, aspect=2.2, draw=orange!70!black, fill=orange!12, line width=0.8pt,
                align=center, inner sep=2pt, font=\small},
  engv/.style ={draw=orange!75!black, fill=orange!14, line width=0.9pt, rounded corners=2pt,
                align=center, inner sep=6pt, font=\small, text width=4.2cm},
  engp/.style ={draw=violet!70!black, fill=violet!12, line width=0.9pt, rounded corners=2pt,
                align=center, inner sep=6pt, font=\small, text width=4.2cm},
  engh/.style ={draw=teal!75!black, fill=teal!12, line width=0.9pt, rounded corners=2pt,
                align=center, inner sep=6pt, font=\small, text width=4.2cm},
  note/.style ={draw=gray!55!black, dashed, fill=yellow!8, line width=0.6pt,
                align=left, inner sep=6pt, font=\footnotesize, text width=15cm},
  fl/.style   ={-{Stealth[length=6pt]}, line width=0.8pt, gray!55!black},
}

\begin{document}
\begin{tikzpicture}[node distance=0.9cm and 1.1cm]

\node[font=\large] (title) {Kiến trúc pipeline trích xuất hóa đơn — dịch vụ OCR};

\node[io, below=0.5cm of title] (input) {Ảnh hóa đơn (JPG/PNG) hoặc PDF quét\\ \footnotesize (\_to\_images: PDF $\rightarrow$ ảnh 200 dpi, ảnh $\rightarrow$ RGB)};
\node[proc, below=of input] (api) {API \texttt{POST /extract}\ \ (fileUrl, model)};
\node[dec, below=of api] (route) {Định tuyến theo\\ tham số \texttt{model}};

% ---- 3 engine song song ----
\node[engv, below=1.4cm of route.south, xshift=-5.6cm] (v)
  {\textbf{model = vietocr}\\[2pt] EasyOCR (CRAFT) phát hiện dòng\\ $\rightarrow$ VietOCR (VGG19 + Transformer) nhận dạng};
\node[engp, below=1.4cm of route.south] (p)
  {\textbf{model = paddleocr}\\[2pt] PaddleOCR PP-OCRv4: DBNet phát hiện $\rightarrow$ phân loại góc $\rightarrow$ SVTR\_LCNet nhận dạng (CTC)};
\node[engh, below=1.4cm of route.south, xshift=5.6cm] (h)
  {\textbf{model = paddledet\_viet}\\[2pt] DBNet (Paddle) phát hiện $\rightarrow$ VietOCR nhận dạng (biến thể lai)};

\node[proc, below=1.5cm of p, text width=8.2cm] (norm)
  {Danh sách dòng chuẩn hóa\\ \texttt{[\{ text, box(x0,y0,x1,y1), prob \}]}};
\node[proc, below=of norm, text width=8.2cm] (parse)
  {\texttt{parse\_invoice()}: khử nghiêng + gom hàng theo vị trí\\ + biểu thức chính quy + chuẩn hoá số/dấu tiếng Việt};
\node[proc, below=of parse, text width=8.2cm] (fields)
  {Trường hóa đơn: MST người bán · Số HĐ · Ngày\\ · Tiền hàng · Tiền thuế · Tổng thanh toán · Tiền tệ};
\node[io, below=of fields, text width=8.2cm] (write)
  {Ghi Invoice + InvoiceItem (đối chiếu Đơn mua PO)};

% ---- luồng ----
\draw[fl] (input) -- (api);
\draw[fl] (api) -- (route);
\draw[fl] (route.south) -- (v.north);
\draw[fl] (route.south) -- (p.north);
\draw[fl] (route.south) -- (h.north);
\draw[fl] (v.south) |- (norm.north);
\draw[fl] (p.south) -- (norm.north);
\draw[fl] (h.south) |- (norm.north);
\draw[fl] (norm) -- (parse);
\draw[fl] (parse) -- (fields);
\draw[fl] (fields) -- (write);

% ---- ghi chú ----
\node[note, below=0.7cm of write] (nt)
  {\textbf{Ghi chú.} Cả ba engine đều tự chủ, chạy offline và fine-tune được trên dữ liệu hóa đơn tự gán nhãn.
   Giao diện cho người dùng chọn \texttt{paddleocr} (mặc định) hoặc \texttt{vietocr};
   \texttt{paddledet\_viet} là biến thể lai (DBNet + VietOCR) dùng để so sánh/đánh giá.
   Backend .NET lưu lựa chọn ở cột \texttt{ocr\_jobs.ocrEngine} và gọi dịch vụ qua trường \texttt{model}.};

\end{tikzpicture}
\end{document}
